% !TEX root = ../main.tex

\chapter*{体例与读法}
\addcontentsline{toc}{chapter}{体例与读法}

\section*{层次}

\begin{enumerate}[label=\textbf{\arabic*.}]
  \item \textbf{【正文·仿罗什风】}——通读主文。
  \item \textbf{【今译意】}——现代汉语，与主文对照。
  \item \textbf{【附注】}——出处、同型删并说明（非每篇皆有长注）。
\end{enumerate}

\noindent 本册\textbf{不}收【底本】与【平行】栏；若需核对文献，请改读 V1／V2 研究译注本。

\section*{信达雅（本册）}

\begin{description}[style=nextline,leftmargin=3em]
  \item[信] 对各篇所据经群／平行之义负责；删并处见附注与书末简表。
  \item[达] 篇章按法义推进，减少无信息增量的重复。
  \item[雅] 仿罗什删冗意译；禁后期大乘／禅宗形上语汇。
\end{description}

\section*{建议读法}

先通读一章正文与今译；遇名相疑问，再以附注中的大正经号查 V1／V2。
篇题为通读新拟，勿与大正经题相混。

\clearpage
